\begin{problem}[Challenge: quasi-compact opens are finite principal unions]
Let $U\subseteq\Spec A$ be open.  Prove that $U$ is quasi-compact if and only if there
exist finitely many $f_1,\dots,f_r\in A$ such that
\[
U=D(f_1)\cup\cdots\cup D(f_r).
\]
\end{problem}

\paragraph{Why this challenge matters.}
It characterizes quasi-compact opens entirely in the principal-open language that will later
support affine gluing and quasi-coherent sheaves.

\paragraph{Useful ideas.}
One direction uses the principal-open basis; the other uses quasi-compactness of each
$D(f_i)$ and stability under finite unions.

\paragraph{Guided hints.}
For every point of $U$, choose a principal open contained in $U$.

\begin{solution}
Assume first that $U$ is quasi-compact.  Because the principal opens form a basis, for each
$\mathfrak p\in U$ there exists $f_{\mathfrak p}\in A$ with
\[
\mathfrak p\in D(f_{\mathfrak p})\subseteq U.
\]
These principal opens cover $U$.  Quasi-compactness gives finitely many points
$\mathfrak p_1,\dots,\mathfrak p_r$ such that
\[
U=D(f_{\mathfrak p_1})\cup\cdots\cup D(f_{\mathfrak p_r}).
\]

Conversely, each $D(f_i)$ is quasi-compact by the principal-open quasi-compactness theorem,
and a finite union of quasi-compact subspaces is quasi-compact.  Hence any finite union of
principal opens is quasi-compact.
\end{solution}

\paragraph{What happened mathematically.}
The principal opens are not merely a basis; they form a basis of quasi-compact opens.

\paragraph{Extensions.}
Show that intersections of quasi-compact opens in $\Spec A$ are quasi-compact.

\paragraph{Provenance.}
New challenge; a structural consequence of the legacy principal-open basis theorem.
